%% Chapter 1 - Introduction
%% Self-contained chapter file. All sections (1.1 Overall project context;
%% 1.2 Individual component context; 1.3 Problem statement; 1.4 Research gap;
%% 1.5 Aim and objectives; 1.6 Research questions; 1.7 Scope, limitations and
%% assumptions; 1.8 Individual contribution; 1.9 Thesis structure) are inlined
%% below. No external section files are required.
\chapter{Introduction}\label{ch:introduction}

\section{Overall Project Context}\label{sec:overall-project-context}

Modern security operations depend on specialised tooling to discover, assess and
remediate faults across networks and applications. As the scale of this work
grows, the manual translation of an operator's intent into safe, well-scoped tool
actions becomes a significant bottleneck, and mistakes made at this point---an
incorrect target, a malformed instruction, or an out-of-scope action---can affect
production infrastructure. Large Language Models (LLMs) are attractive for this
translation step because they can interpret natural language and emit structured,
machine-readable instructions that drive security tooling
\parencite{deng2024pentestgpt}. Yet a model is a stochastic component: its output
may be malformed, misclassified, hallucinated, or the product of injected
instructions, and a security pipeline that acts on such output without further
checking transfers every one of these errors into the actions it executes.

NeuroShell is an autonomous, natural-language-driven security orchestration
framework developed as a four-component research chain: the Intent Recognition
Engine (IRE), a Dynamic Planner, an Adaptive Execution and Error Recovery engine,
and an AI-driven Vulnerability Analysis engine (project reference R26-CS-018).
The framework is designed to let an analyst express offline security tasks in
natural language and have the resulting validated, structured contracts drive
downstream automation, within an explicitly configured engagement scope. It
therefore sits at the intersection of cybersecurity automation, software
engineering and applied machine learning. Its central concern is the first step
of that chain: the point at which human intent is converted into a
machine-executable form, because an error introduced there propagates through
every stage that follows and must be detected before it takes effect.

This dissertation concentrates on a single component of the framework and asks
whether, in an offline CPU-only deployment, a small instruction-tuned language
model can be made sufficiently reliable to serve as the head of an automated
security pipeline when its raw output is constrained by deterministic
validation.

\section{Individual Component Context}\label{sec:individual-component-context}

The Intent Recognition Engine (IRE) is the first component of the chain
(Component 01). It acts as a domain-specific Natural Language Understanding (NLU)
gateway: it receives natural-language security commands---for example a request
to perform a targeted network discovery scan---and returns a structured,
validated intent contract that downstream components consume without
re-interpreting prose.

The component boundary is deliberately narrow. The IRE is concerned with
understanding and validating \emph{intent}; it does not decide how a task is
scheduled, across which hosts it runs, or how execution failures are recovered.
Those responsibilities belong to the Dynamic Planner (Component 02) and to the
Adaptive Execution and Error Recovery engine (Component 03), which consume the
IRE's output. The handoff is designed to be safe by construction: the output
contract carries an explicit intent type, a resolved target, the recognised
options and any validation findings, and the downstream contract requires that
actions derived for rejected intents must not be executed.

Because the IRE sits at the head of the chain, its reliability bounds the
reliability of the components that follow. A target that is mislabeled, an
intent that is misclassified, or a hallucinated identifier that is accepted would
be inherited by every downstream action. The component therefore couples an
instruction-tuned language model with deterministic validation and guard stages
whose purpose is to contain model error before it can propagate to a planner or
executor. This dissertation evaluates whether that guarded design achieves
sufficiently accurate and error-resilient intent recognition to justify acting on
model-produced security instructions in a controlled, offline setting.

\section{Problem Statement}\label{sec:problem-statement}

Security operations are bottlenecked by the manual translation of operator intent
into tool actions. A penetration tester or analyst must map free-form
instructions---``scan the internal network for EternalBlue'', ``check this host
for CVE-2021-44228''---onto commands, targets and constraints while respecting
authorisation boundaries. Recent work shows that Large Language Models (LLMs) can
act as capable copilots for parts of this work, producing realistic attack plans
and driving tools in controlled environments \parencite{deng2024pentestgpt,
happe2023pwned}. Using an LLM for this translation step is attractive precisely
because natural language is the operator's native interface.

The problem is that a language model is not a reliable translator. Its output can
be malformed, misclassified, or wholly fabricated: LLMs hallucinate identifiers
and parameters \parencite{ji2023hallucination}, and they can be diverted by
adversarially crafted utterances \parencite{greshake2023indirect,
liu2023promptinjection}. In a purely conversational copilot these errors are
display text that a human can vet. In an \emph{automated} security pipeline,
where intent is converted directly into actions, each such error becomes a
candidate machine-dispatched action: a fabricated CVE is looked up, a wrong
target is scanned, or an injected instruction is obeyed. Acting on unvalidated
model output therefore transfers every residual error of the model into the
security toolchain it drives.

The intent-recognition step concentrates this risk. A copilot that pipelines raw
LLM output into a planner and executor must, at minimum, emit a well-formed,
schema-conformant, authorised intent structure before any downstream component
acts on it. Unless that structured output is validated deterministically, the
entire framework inherits the model's failure modes. This dissertation addresses
the resulting problem: \emph{how can a small, offline, CPU-only language model be
made sufficiently reliable to serve as the intent-recognition head of an
automated security pipeline, when its raw output is constrained by deterministic
validation rather than by a human reviewer or a cloud back-end?}

\section{Research Gap}\label{sec:research-gap}

Three related gaps motivate this work.

First, LLM copilots for security tooling validate their output informally, if at
all. PENTESTGPT decomposes penetration-testing tasks across a reasoning,
generation and guidance chain and does not subject its output to deterministic schema, scope
or authorisation validation before machine consumption
\parencite{deng2024pentestgpt}; the sparring-partner systems demonstrated by
Happe and Cito produce plausible next actions consumed by a human
\parencite{happe2023pwned}. When a copilot's output is fed to an
\emph{automated} toolchain---the natural integration point for intent
recognition---the requirement changes from plausibility to correctness: a
hallucinated CVE, malformed structure, or injected text becomes a
machine-dispatched action rather than display text. The OWASP Top 10 for LLM
applications catalogues prompt injection, insecure output handling and related
risks as first-class application concerns \parencite{owasp2025}, yet the
systematic defence of machine-consumed intent remains largely unaddressed in the
surveyed copilot work.

Second, the defence literature either assumes interactive, cloud-hosted
generation or takes the retrieval-augmented route. Sampling-based methods for
detecting fabricated output, such as SelfCheckGPT \parencite{manakul2023selfcheckgpt},
semantic uncertainty \parencite{kuhn2023semantic} and semantic entropy
\parencite{farquhar2024semantic}, and the self-consistency decoding strategy
\parencite{wang2023selfconsistency}, all assume repeated or interactive
generation with generous compute. Retrieval-augmented generation
\parencite{lewis2020rag} grounds answers externally, but introduces retrieved
data as a prompt-injection surface \parencite{greshake2023indirect} and requires
corpora and embeddings infrastructure at inference time. The operational setting
targeted here---a small offline model running unattended on an operator's
laptop, with latency budget that rules out a retry loop or a verifier LLM---is
therefore not covered by direct prior measurements.

Third, to the best of the review conducted here, no existing
autonomous penetration-testing framework ships a dedicated, isolated
natural-language understanding layer that reliably converts free-form
offensive-security instructions into validated, schema-constrained machine contracts; tools either accept structured input or
rely on general-purpose LLMs without domain-specific validation. Leaving this
step to a general chat model without a purpose-built validation layer is exactly
the unsafe pattern identified above. The research gap this dissertation fills is
therefore the design and \emph{measurement} of a guarded, retrieval-free intent
recognition engine---an isolated NLU component with deterministic validation---
for offline, single-device security deployments.

\section{Aim and Objectives}\label{sec:aim-and-objectives}

The aim of this dissertation is to design, implement and evaluate a guarded
intent-recognition engine that constrains the raw output of a small, offline,
instruction-tuned language model with deterministic validation, so that the
combined system can act as a reliable head of an automated security pipeline in a
CPU-only, single-device deployment. To achieve this aim the following objectives
are pursued:

\begin{enumerate}
\item To design and implement the Intent Recognition Engine (IRE): a component
that receives free-text security requests and emits strictly
schema-conformant, authorised intent contracts, consuming a small offline
language model as a proposition generator whose output is validated
deterministically before it is handed onward.
\item To construct a purpose-built evaluation data set of security prompts,
covering well-formed requests and deliberately flawed prompts across the flaw
categories that matter in practice (adversarial phrasing, out-of-scope targets,
ambiguous or contradictory requests, fabricated parameters and identifiers, and
role violations), each annotated with gold labels.
\item To establish an evaluation protocol that quantifies, on the committed
experiment record, how much the guarded pipeline improves structured intent
recognition and flaw catching over raw model output, and that abates each guard
layer independently to estimate its individual contribution.
\item To compare two offline language-model backends under the identical
guardrail stack, in order to measure how much of the achieved safety is
backend-independent rather than an artefact of one model.
\end{enumerate}

Each objective is addressed by a corresponding chapter of this dissertation:
the design is developed in Chapter~\ref{ch:product-design}, the data set and protocol in
Chapter~\ref{ch:methodology}, the implementation in
Chapter~\ref{ch:implementation-testing}, and the evaluation and backend comparison in
Chapter~\ref{ch:results}.

\section{Research Questions}\label{sec:research-questions}

The central question guiding this dissertation is whether a small, offline,
unattended language model can be made reliable enough to drive an automated
security copilot once its output is subject to deterministic validation.
From this central question three concrete research questions are derived:

\begin{enumerate}
\item \textbf{RQ1 (pipeline value).} How much does a stacked validation
pipeline improve structured intent-recognition accuracy and flaw catching over
raw language-model output?
\item \textbf{RQ2 (layer contribution).} Which individual guard layer
contributes the most to the flaw-catch rate, and which layers interact
destructively when combined?
\item \textbf{RQ3 (generalisability).} Do the guardrails generalise across two
different offline language-model backends, or is the measured safety an
artefact of a single model?
\end{enumerate}

The study is conducted as an empirical, measurement-based investigation rather
than a hypothesis-testing one: the outcomes of RQ1--RQ3 are quantified from the
committed experiment record reported in Chapter~\ref{ch:results}, and no
statistical significance claim is attached to the single-run measurements
(see Section~\ref{sec:scope-limitations-assumptions}).

\section{Scope, Limitations and Assumptions}\label{sec:scope-limitations-assumptions}

\subsection*{Scope}

This dissertation is an individual study of a single component, the Intent
Recognition Engine (Component 01 of the NeuroShell framework, project reference
R26-CS-018), and does not implement the downstream planner, executor or
vulnerability-analysis components. Within that boundary, the scope covers:
(1) a guarded NLU pipeline that accepts free-text security requests and emits
validated intent contracts; (2) the deterministic validation and security
countermeasures that constrain the model's output; (3) an offline, CPU-only,
single-device deployment in which nothing is retrieved at inference time and no
request leaves the device; and (4) a closed intent taxonomy and an
operator-configured engagement scope that bounds permitted targets and
techniques. Evaluation is conducted on a curated, annotated data set of prompts
in a lab harness, not against live production traffic.

\subsection*{Limitations}

The principal limitation is evidential: the results reported in this
dissertation rest on a bounded, committed experiment record of
single-run point estimates, without confidence intervals or seed sweeps, and
therefore carry residual uncertainty that is discussed in
Chapter~\ref{ch:discussion}. Only two offline language-model backends are
compared under the guardrail stack, and the evaluation data set is
author-authored rather than drawn from observed operational traffic.
Measurements are specific to the evaluation hardware and CPU-only inference
runtime, so latency and throughput figures are indicative rather than
generalisable. The framework-level benefit of the IRE is validated at the
component level; end-to-end orchestration with the downstream components is
outside the measured evidence presented here.

\subsection*{Assumptions}

The following assumptions underpin the design and evaluation:
\begin{itemize}
\item The operator configures the engagement scope and authorisation policy
correctly; the scope guard can only enforce the boundary it is given.
\item The deployment device and the downloaded model weights are trusted;
physical compromise of the device and backdoors embedded in the model weights
are out of scope.
\item Downstream components consume the intent contract as specified and
honour the contract rule that actions derived for rejected intents must not be
executed.
\item Because the pipeline performs no retrieval and renders no free-form
output, indirect prompt injection via retrieved content and tool injection via
unfiltered rendered text are structurally out of scope, narrowing the attack
surface to the classification of the single user utterance.
\item The annotated prompts in the evaluation data set are representative of
the phrasing an operator would use with the system.
\end{itemize}

\section{Individual Contribution}\label{sec:individual-contribution}

This dissertation is an individual thesis on Component 01 of the NeuroShell
four-component research chain. The candidate's individual contribution is the
design, implementation and evaluation of the Intent Recognition Engine, and
comprises: (1) the architecture of the guarded NLU gateway, in which a small
offline language model acts as a proposition generator and deterministic stages
validate structure, target syntax and authorisation before an intent is
emitted; (2) the security countermeasures layered around the core---including
the authorisation checks, the adversarial-utterance detection, and the
self-consistency and hallucination-validation gate---that enforce the safety
invariant that every accepted intent is schema-valid, authorised and in scope;
(3) the authoring of the annotated evaluation data set covering well-formed and
deliberately flawed prompt categories; (4) the evaluation harness and the
per-layer ablation protocol used to quantify the pipeline's value and each
layer's contribution; and (5) the empirical comparison of two offline
language-model backends under the identical guardrail stack.

The other components of the NeuroShell chain---the dynamic planner, the
adaptive executor, and the vulnerability-analysis engine---provide the context
in which the IRE operates and are described only to the extent needed to define
the IRE's interface contract. The candidate takes responsibility for the
correctness of the IRE implementation and for every result reported in this
dissertation.

\section{Thesis Structure}\label{sec:thesis-structure}

This dissertation is organised as follows. \emph{Chapter~\ref{ch:introduction}}
(the present chapter) establishes the problem, the research gap, the aim,
objectives, research questions and scope of the study.
\emph{Chapter~\ref{ch:literature-review}} reviews related work on LLM copilots
for security, adversarial attacks on LLM-integrated applications, and
hallucination-detection and uncertainty methods, and situates the guarded
intent-recognition approach within it.
\emph{Chapter~\ref{ch:methodology}} presents the research approach: the
objectives-to-method mapping, the design of the evaluation data set, the
metrics, and the experiment and ablation protocol used to answer RQ1--RQ3.
\emph{Chapter~\ref{ch:product-design}} develops the product and solution
design of the Intent Recognition Engine, including the pipeline stages, the
security countermeasures and the intent-contract interface.
\emph{Chapter~\ref{ch:implementation-testing}} describes the implementation,
the integration of the two offline backends, and the automated test suite.
\emph{Chapter~\ref{ch:results}} reports the results of the pipeline-value,
layer-ablation and backend-comparison experiments and interprets them against
the research questions.
\emph{Chapter~\ref{ch:discussion}} discusses the findings, the threats to
validity, and the residual failure modes of the guarded design.
\emph{Chapter~\ref{ch:commercialisation}} considers the commercialisation and
sustainability of the approach.
\emph{Chapter~\ref{ch:conclusion}} concludes the dissertation, restates the
contributions, and outlines directions for future work. Supplementary material
appears in the appendices: the data-set and evaluation instruments
(Appendix~\ref{app:a}) and the declaration on academic integrity and the use of
artificial-intelligence tools (Appendix~\ref{app:b}).